\section{A paid component: constructive partitioning and the boundary contribution}\label{sec:boundary}
Unlike cancellation of the full scalar, the boundary term can be controlled by support size while preserving the actual objects\cite{archivephysical}.
\begin{theorem}[Collective bound for the actual weighted boundary]\label{thm:boundary}
There is a nonnegative smooth partition $\sum_j\eta_{j,H}=1$ with finite overlap, each support of length $H$, and
$H^r\|\eta_{j,H}^{(r)}\|_\infty\le C_r$ uniformly in $j,H$.
Classify labels only inside the physical set $\mathcal J_x$ of supports meeting $[x/2,x]$:
\[
\mathcal J_\partial=\{j\in\mathcal J_x:\operatorname{supp}\eta_{j,H}\not\subset(x/2,x)\}.
\]
Let $\theta_\partial=\sum_{j\in\mathcal J_\partial}\eta_{j,H}$ be their total weight and set $\theta_{\rm int}=1-\theta_\partial$. Then
\[
\left|\Cx-\ip{\theta_{\rm int}w}{A(\theta_{\rm int}\beta)}\right|
\ll xQ^2x^{-11/64+o(1)}.
\]
\end{theorem}
\begin{proof}
Take $\chi(s)=\exp[-1/(1-4s^2)]$ for $|s|<1/2$ and zero elsewhere, and set
\[
D(s)=\sum_{j\in\ZZ}\chi(s-j/2),\qquad
\eta_{j,H}(t)=\frac{\chi(t/H-j/2)}{D(t/H)}.
\]
The function $D$ is positive, smooth and $1/2$-periodic: every point lies within $1/4$ of a half-integer. All derivatives of $D$ and $D^{-1}$ are bounded on one period. The chain rule gives the required derivative estimates.
Each closed support is contained in $[(j-1)H/2,(j+1)H/2]$, at most three closed supports overlap, and $|\mathcal J_x|=O(x/H+1)$.
Since every classified support meets the physical closed interval, a boundary support must contain at least one endpoint. There are therefore at most six such supports, and
\[
\operatorname{supp}_{I_x}\theta_\partial
\subset [x/2,x/2+H]\cup[x-H,x],
\quad |\operatorname{supp}_{I_x}\theta_\partial|=O(H+1).
\]
Only pointwise envelopes for the actual coefficients are needed. The divisor bound $\tau(n)\ll_\epsilon n^\epsilon$ gives $|\beta(n)|\le1+\tau(n)$.
By \eqref{eq:hybrid}, $0\le b_x^{(z)}(n)\le W_z2^{\omega(n)}\le W_z\tau(n)$, so
$|\beta(n)|+|w(n)|\le x^{o(1)}$ uniformly on the actual interval, for fixed $K$.
The divisor bound follows prime power by prime power: for $p\ge2^{1/\epsilon}$,
$k+1\le2^k\le p^{\epsilon k}$; for each of the finitely many remaining primes,
$\sup_k(k+1)p^{-\epsilon k}$ is finite.

Thus each boundary vector has norm $H^{1/2}x^{o(1)}$, while each full vector has norm $x^{1/2+o(1)}$.
Expand exactly, counting the double-boundary term only once:
\[
\Cx-\ip{w_{\rm int}}{A\beta_{\rm int}}
=\ip{w_\partial}{A\beta}+\ip{w_{\rm int}}{A\beta_\partial}.
\]
Lemma~\ref{lem:upper}, namely $\norm A\ll Q^2$, now gives
$Q^2\sqrt{xH}x^{o(1)}=xQ^2x^{-11/64+o(1)}$.
\end{proof}
The theorem bounds the \emph{collective contribution} of all weighted ordered pairs with a boundary label, not the absolute value of each positive polarization package.
It makes an admissible explicit choice of partition; it does not assert derivative bounds for every previously specified partition.
The interior cutoff is genuinely smooth, but this does not turn $\Lambda$, the M\"obius coefficients, or the hybrid main term into smooth coefficients.
If an interior saving $\delta_{\rm int}$ is proved in the future, this recombination certifies a full saving of $\min(\delta_{\rm int},11/64)$. The signed interior estimate remains unresolved here.
